\documentclass[12pt,a4paper]{article}
\usepackage[UTF8]{ctex}
\usepackage{amsmath,amssymb,amsthm}
\usepackage{geometry}

\geometry{margin=2.5cm}

\title{导纳控制中忽略虚拟质量$M$的优势和劣势}
\author{第11章补充说明}
\date{}

\begin{document}

\maketitle

\section{问题提出}

\textbf{问题}：在导纳控制中，忽略虚拟质量$M$（设置$M = 0$）有什么优势和劣势？

\textbf{导纳控制公式}（标准形式）：
\begin{equation}
M\ddot{x}_d + B\dot{x} + Kx = f_{\text{ext}} \tag{11.64}
\end{equation}

\textbf{忽略$M$后的公式}：
\begin{equation}
B\dot{x} + Kx = f_{\text{ext}} \tag{11.64-M=0}
\end{equation}

\textbf{求解期望加速度}（当$M \neq 0$时）：
\begin{equation}
\ddot{x}_d = M^{-1}(f_{\text{ext}} - B\dot{x} - Kx) \tag{11.66}
\end{equation}

\textbf{忽略$M$后}（当$M = 0$时）：
\begin{equation}
\dot{x}_d = B^{-1}(f_{\text{ext}} - Kx) \tag{11.66-M=0}
\end{equation}

注意：当$M = 0$时，我们直接计算期望速度$\dot{x}_d$，而不是加速度$\ddot{x}_d$。

\section{忽略$M$的优势}

\subsection{优势1：不需要测量或计算加速度}

\textbf{当$M \neq 0$时}：
\begin{itemize}
\item 需要计算期望加速度$\ddot{x}_d = M^{-1}(f_{\text{ext}} - B\dot{x} - Kx)$
\item 需要测量或估计当前加速度$\ddot{x}$（用于某些控制律）
\item 加速度测量通常是噪声的
\end{itemize}

\textbf{当$M = 0$时}：
\begin{itemize}
\item 直接计算期望速度$\dot{x}_d = B^{-1}(f_{\text{ext}} - Kx)$
\item 不需要加速度信息
\item 只需要位置$x$和速度$\dot{x}$（或外力$f_{\text{ext}}$）
\end{itemize}

\textbf{优势}：
\begin{itemize}
\item 减少传感器需求（不需要加速度计）
\item 避免加速度测量的噪声问题
\item 计算更简单
\end{itemize}

\subsection{优势2：避免质量补偿的矛盾}

\textbf{问题描述}：

在阻抗控制中，如果使用质量补偿项$\tilde{\Lambda}(\theta)\ddot{x}$，存在矛盾：
\begin{enumerate}
\item 测量加速度$\ddot{x}$
\item 计算质量补偿项$\tilde{\Lambda}(\theta)\ddot{x}$
\item 应用控制力矩
\item 加速度$\ddot{x}$改变（因为力矩改变了）
\item 需要重新测量，但测量值已经过时
\end{enumerate}

\textbf{在导纳控制中}：

虽然导纳控制不直接使用质量补偿项，但如果$M \neq 0$，仍然需要计算加速度$\ddot{x}_d$，这可能涉及类似的问题。

\textbf{当$M = 0$时}：
\begin{itemize}
\item 不需要计算加速度
\item 避免了质量补偿的矛盾
\item 控制律更直接
\end{itemize}

\subsection{优势3：简化控制律}

\textbf{完整导纳控制律}（$M \neq 0$）：

\textbf{步骤1}：计算期望加速度
\begin{equation}
\ddot{x}_d = M^{-1}(f_{\text{ext}} - B\dot{x} - Kx)
\end{equation}

\textbf{步骤2}：转换为关节加速度
\begin{equation}
\ddot{\theta}_d = J^{\dagger}(\theta)(\ddot{x}_d - \dot{J}(\theta)\dot{\theta})
\end{equation}

\textbf{步骤3}：计算关节力矩
\begin{equation}
\tau = M(\theta)\ddot{\theta}_d + h(\theta, \dot{\theta})
\end{equation}

\textbf{简化导纳控制律}（$M = 0$）：

\textbf{步骤1}：计算期望速度
\begin{equation}
\dot{x}_d = B^{-1}(f_{\text{ext}} - Kx)
\end{equation}

\textbf{步骤2}：转换为关节速度
\begin{equation}
\dot{\theta}_d = J^{\dagger}(\theta)\dot{x}_d
\end{equation}

\textbf{步骤3}：控制关节速度
\begin{equation}
\dot{\theta} = \dot{\theta}_d
\end{equation}

\textbf{优势}：
\begin{itemize}
\item 不需要计算加速度
\item 不需要逆动力学（如果使用速度控制）
\item 不需要$\dot{J}(\theta)$（雅可比的时间导数）
\item 实现更简单
\end{itemize}

\subsection{优势4：更适合速度控制执行器}

\textbf{许多机器人系统}：
\begin{itemize}
\item 只能控制关节速度，不能直接控制关节力矩
\item 例如：步进电机、某些伺服电机
\end{itemize}

\textbf{当$M \neq 0$时}：
\begin{itemize}
\item 需要计算加速度$\ddot{x}_d$
\item 需要积分得到速度$\dot{x}_d = \int \ddot{x}_d dt$
\item 积分可能引入误差
\end{itemize}

\textbf{当$M = 0$时}：
\begin{itemize}
\item 直接计算速度$\dot{x}_d$
\item 不需要积分
\item 更适合速度控制执行器
\end{itemize}

\subsection{优势5：减少计算复杂度}

\textbf{计算复杂度对比}：

\textbf{当$M \neq 0$时}：
\begin{itemize}
\item 需要计算$M^{-1}$（矩阵求逆）
\item 需要计算$\ddot{x}_d$
\item 需要计算$\dot{J}(\theta)$（雅可比的时间导数）
\item 需要逆动力学计算
\end{itemize}

\textbf{当$M = 0$时}：
\begin{itemize}
\item 只需要计算$B^{-1}$（如果$B$是对角矩阵，求逆很简单）
\item 不需要计算加速度
\item 不需要$\dot{J}(\theta)$
\item 不需要逆动力学（如果使用速度控制）
\end{itemize}

\textbf{优势}：
\begin{itemize}
\item 计算更快
\item 实时性更好
\item 适合计算资源有限的系统
\end{itemize}

\subsection{优势6：避免数值问题}

\textbf{当$M$很小时}：
\begin{itemize}
\item $M^{-1}$可能很大，导致数值不稳定
\item 加速度$\ddot{x}_d = M^{-1}(f_{\text{ext}} - B\dot{x} - Kx)$可能很大
\item 可能导致控制力矩过大
\end{itemize}

\textbf{当$M = 0$时}：
\begin{itemize}
\item 避免了$M^{-1}$的计算
\item 避免了数值不稳定问题
\item 控制更稳定
\end{itemize}

\section{忽略$M$的劣势}

\subsection{劣势1：不能模拟质量特性}

\textbf{物理意义}：

虚拟质量$M$表示机器人应该表现的"惯性"特性。如果$M = 0$，机器人不能模拟质量。

\textbf{影响}：
\begin{itemize}
\item 用户感觉不到"重量"
\item 机器人表现得像没有质量
\item 不能模拟真实物体的惯性特性
\end{itemize}

\textbf{例子}：

\textbf{场景}：用户推动机器人末端执行器

\textbf{当$M \neq 0$时}：
\begin{itemize}
\item 用户需要施加力来加速机器人
\item 感觉像推动一个有质量的物体
\item 符合物理直觉
\end{itemize}

\textbf{当$M = 0$时}：
\begin{itemize}
\item 用户推动时，机器人立即响应
\item 感觉像推动一个没有质量的物体
\item 不符合物理直觉
\end{itemize}

\subsection{劣势2：响应可能过快}

\textbf{问题描述}：

当$M = 0$时，系统响应可能过快，导致：
\begin{itemize}
\item 不自然的运动
\item 可能产生振荡
\item 用户感觉不舒服
\end{itemize}

\textbf{数学分析}：

\textbf{当$M \neq 0$时}（二阶系统）：
\begin{equation}
M\ddot{x}_d + B\dot{x} + Kx = f_{\text{ext}}
\end{equation}

传递函数：
\begin{equation}
\frac{X_d(s)}{F_{\text{ext}}(s)} = \frac{1}{Ms^2 + Bs + K}
\end{equation}

\textbf{当$M = 0$时}（一阶系统）：
\begin{equation}
B\dot{x} + Kx = f_{\text{ext}}
\end{equation}

传递函数：
\begin{equation}
\frac{X_d(s)}{F_{\text{ext}}(s)} = \frac{1}{Bs + K}
\end{equation}

\textbf{对比}：
\begin{itemize}
\item 二阶系统有自然频率和阻尼比，可以设计响应特性
\item 一阶系统响应更快，但可能不够平滑
\end{itemize}

\subsection{劣势3：不能提供惯性阻尼}

\textbf{物理意义}：

质量$M$提供惯性，可以"平滑"运动，减少突然的变化。

\textbf{当$M \neq 0$时}：
\begin{itemize}
\item 质量$M$提供惯性阻尼
\item 运动更平滑
\item 减少振荡
\end{itemize}

\textbf{当$M = 0$时}：
\begin{itemize}
\item 没有惯性阻尼
\item 只能依靠阻尼$B$来减少振荡
\item 如果$B$不够大，可能产生振荡
\end{itemize}

\subsection{劣势4：不适合模拟低质量环境}

\textbf{问题描述}：

如果要模拟低质量环境（例如，模拟轻量级物体），需要小的$M$，但不能设置为$M = 0$。

\textbf{当$M = 0$时}：
\begin{itemize}
\item 完全不能模拟质量
\item 不能区分"轻"和"无质量"
\item 限制了应用场景
\end{itemize}

\textbf{当$M \neq 0$但很小时}：
\begin{itemize}
\item 可以模拟轻量级物体
\item 仍然有惯性特性
\item 更灵活
\end{itemize}

\subsection{劣势5：稳态误差可能不同}

\textbf{当$M \neq 0$时}：

对于阶跃外力$f_{\text{ext}} = f_0$（常数），稳态时$\ddot{x}_d = 0$，$\dot{x} = 0$：
\begin{equation}
Kx_{ss} = f_0 \quad \Rightarrow \quad x_{ss} = \frac{f_0}{K}
\end{equation}

\textbf{当$M = 0$时}：

对于阶跃外力$f_{\text{ext}} = f_0$（常数），稳态时$\dot{x}_d = 0$：
\begin{equation}
Kx_{ss} = f_0 \quad \Rightarrow \quad x_{ss} = \frac{f_0}{K}
\end{equation}

\textbf{结论}：稳态误差相同，但瞬态响应不同。

\subsection{劣势6：不能利用质量的稳定性}

\textbf{物理意义}：

质量$M$可以提供稳定性，特别是在高频扰动下。

\textbf{当$M \neq 0$时}：
\begin{itemize}
\item 质量$M$可以"吸收"高频扰动
\item 系统更稳定
\item 对噪声不敏感
\end{itemize}

\textbf{当$M = 0$时}：
\begin{itemize}
\item 没有质量来"吸收"扰动
\item 对噪声更敏感
\item 可能需要更大的阻尼$B$来稳定系统
\end{itemize}

\section{实际应用建议}

\subsection{何时忽略$M$？}

\textbf{适合忽略$M$的情况}：
\begin{enumerate}
\item \textbf{速度控制执行器}：
\begin{itemize}
\item 只能控制速度，不能控制力矩
\item 设置$M = 0$，直接计算速度
\end{itemize}

\item \textbf{计算资源有限}：
\begin{itemize}
\item 需要快速计算
\item 不能承受复杂的逆动力学计算
\end{itemize}

\item \textbf{不需要模拟质量}：
\begin{itemize}
\item 任务不要求模拟质量特性
\item 只需要柔顺性（弹簧-阻尼器）
\end{itemize}

\item \textbf{加速度测量不可用}：
\begin{itemize}
\item 没有加速度计
\item 加速度测量噪声太大
\end{itemize}

\item \textbf{轻量级机器人}：
\begin{itemize}
\item 机器人本身很轻
\item 质量对用户不明显
\item 可以忽略虚拟质量
\end{itemize}
\end{enumerate}

\subsection{何时保留$M$？}

\textbf{适合保留$M$的情况}：
\begin{enumerate}
\item \textbf{需要模拟质量特性}：
\begin{itemize}
\item 任务要求模拟真实物体的惯性
\item 用户需要感觉到"重量"
\end{itemize}

\item \textbf{需要平滑响应}：
\begin{itemize}
\item 需要利用质量的惯性阻尼
\item 减少振荡和不平滑的运动
\end{itemize}

\item \textbf{力矩控制执行器}：
\begin{itemize}
\item 可以直接控制关节力矩
\item 可以计算加速度并转换为力矩
\end{itemize}

\item \textbf{需要精确控制}：
\begin{itemize}
\item 需要完整的二阶系统特性
\item 需要设计自然频率和阻尼比
\end{itemize}

\item \textbf{模拟低质量环境}：
\begin{itemize}
\item 需要模拟轻量级物体（$M$很小但不为零）
\item 仍然需要惯性特性
\end{itemize}
\end{enumerate}

\section{折中方案}

\subsection{使用很小的$M$}

\textbf{方法}：不设置$M = 0$，而是设置$M$为一个很小的值。

\textbf{优势}：
\begin{itemize}
\item 保留了质量特性
\item 避免了$M^{-1}$的数值问题（如果$M$不是太小）
\item 仍然可以简化某些计算
\end{itemize}

\textbf{劣势}：
\begin{itemize}
\item 仍然需要计算加速度
\item 仍然需要$M^{-1}$
\end{itemize}

\subsection{使用对角$M$}

\textbf{方法}：使用对角质量矩阵$M = \text{diag}(m_1, m_2, \ldots, m_n)$。

\textbf{优势}：
\begin{itemize}
\item $M^{-1}$的计算很简单（只需对每个元素求倒数）
\item 减少了计算复杂度
\item 仍然可以模拟质量
\end{itemize}

\subsection{自适应$M$}

\textbf{方法}：根据任务需求动态调整$M$。

\textbf{策略}：
\begin{itemize}
\item 在需要快速响应时，减小$M$
\item 在需要平滑运动时，增大$M$
\item 在不需要质量特性时，设置$M = 0$
\end{itemize}

\section{总结}

\subsection{优势总结}

\begin{enumerate}
\item \textbf{不需要加速度信息}：减少传感器需求，避免噪声
\item \textbf{避免质量补偿矛盾}：简化控制律
\item \textbf{更适合速度控制}：直接计算速度，不需要积分
\item \textbf{计算更简单}：不需要逆动力学，不需要$\dot{J}(\theta)$
\item \textbf{避免数值问题}：避免$M^{-1}$的数值不稳定
\end{enumerate}

\subsection{劣势总结}

\begin{enumerate}
\item \textbf{不能模拟质量}：用户感觉不到"重量"
\item \textbf{响应可能过快}：可能不自然，可能振荡
\item \textbf{没有惯性阻尼}：只能依靠阻尼$B$
\item \textbf{不适合模拟低质量环境}：不能区分"轻"和"无质量"
\item \textbf{对噪声更敏感}：没有质量来"吸收"扰动
\end{enumerate}

\subsection{选择建议}

\textbf{忽略$M$（$M = 0$）}：
\begin{itemize}
\item 速度控制执行器
\item 计算资源有限
\item 不需要模拟质量
\item 轻量级机器人
\end{itemize}

\textbf{保留$M$（$M \neq 0$）}：
\begin{itemize}
\item 需要模拟质量特性
\item 需要平滑响应
\item 力矩控制执行器
\item 需要精确控制
\end{itemize}

\textbf{折中方案}：
\begin{itemize}
\item 使用很小的$M$
\item 使用对角$M$
\item 自适应$M$
\end{itemize}

\end{document}

